% SPDX-License-Identifier: Apache-2.0
\section{Who This Chip Is, and Why It Restarted}
\label{sec:x-syscon}

Three questions a program wants answered by the hardware and by
nothing else. Which design am I running on. Why did I start. And how
do I hand a word to whatever runs next.

\code{Syscon} answers all three in seven words. Four are constants
from the type's parameters: an identifier, a version, and two words of
build stamp, which a build can fill with a commit and a date so that a
program can say which bitstream it is running on rather than which one
it was compiled for. The other three are a cause, a reset and a
scratch.

\subsection{The Cause Accumulates}

\code{cause} has a bit per reason, and they are kept until a program
clears them by writing ones. It is tempting to make such a register
hold only the last reason, and that is the wrong choice: the
interesting case is the one where two things went wrong, and a
register that overwrote itself would lose exactly that case. The
example ends by pressing the button while the software bit is still
set, and reads back both.

The four reasons are the power coming on, the button, the watchdog
and a program asking. The first three arrive as wires, so a board
decides what they mean; the fourth is this block's own doing.

\subsection{A Key, and Why}

A write to \code{reset} restarts the system only if it carries
\code{KEY}. A wild store from a program that has lost its way should
not restart the machine, and the cheapest defence is to make the
value improbable. The example writes a zero and checks that nothing
happened.

The reset is held for eight cycles rather than one, because a reset
that lasted a single cycle would be a pulse that some parts of a
design saw and others did not.

\subsection{What A Reset Does Not Touch}

\code{scratch} and \code{cause} are the two registers a reset leaves
alone, and that is what makes them worth having. A bootloader leaves a
word in \code{scratch} and restarts into the program that reads it; a
program reads \code{cause} to find out whether it is starting for the
first time or picking up after a watchdog.

The block does not reset itself, which is how those two survive. That
is a decision rather than an oversight, and it is the reason this part
sits outside whatever the reset line drives.

\begin{figure}[htbp]
\centering
\input{syscon_timing.tex}
\caption{A reset the program asked for: the key written, the line
held for eight cycles while \code{hold} counts down, the cause coming
back with the software bit, and the scratch word untouched through
all of it.}
\label{fig:x-syscon}
\end{figure}

\lstinputlisting[caption={\code{ex\_syscon.rs}. Run with \code{bazel run //lib/examples:ex\_syscon}.},label={lst:x-syscon}]{lib/examples/ex_syscon.rs}

\lstinputlisting[style=ascii,caption={What \code{ex\_syscon} printed when the build ran it: the identity, the build stamp, the cause at each step, and then the block's Verilog.},label={lst:x-syscon-out}]{out_syscon.txt}
